This work investigated the ability of a multilayer perceptron (MLP) to approximate the dynamics of several physical systems with increasing levels of complexity, ranging from the linear harmonic oscillator to the chaotic double pendulum. A common neural network architecture was trained using only time as input, allowing the influence of the underlying system dynamics on the prediction accuracy to be analysed independently of architectural changes.

The results show that the neural network successfully learned the behaviour of all considered systems. For the harmonic oscillator and the damped harmonic oscillator, prediction errors remained extremely low across all experiments, confirming that smooth and regular trajectories can be accurately approximated by a relatively simple feedforward network \cite{cybenko1989approximation}. The forced harmonic oscillator introduced a more demanding learning task due to the interaction between the natural oscillation and the external forcing frequency. Although prediction accuracy remained high for most forcing frequencies, the experiments revealed that high-frequency oscillations produced a noticeable increase in prediction error, suggesting that rapidly varying trajectories are intrinsically more difficult for the network to approximate.

The Duffing oscillator demonstrated that the presence of nonlinear terms alone does not necessarily reduce predictive performance. Across the explored parameter ranges, variations in the nonlinear coefficients produced only minor changes in the approximation error. This indicates that the complexity of the generated trajectory is more relevant than the mathematical form of the governing equation itself.

The simple pendulum produced similar conclusions. Changes in gravity and pendulum length affected prediction accuracy primarily because they modified the oscillation frequency, while variations in the initial angle or initial angular velocity had comparatively little influence. Furthermore, the duration experiments showed that the network could accurately predict long trajectories only up to a certain time horizon, after which the accumulated approximation error became increasingly visible.

The double pendulum represented the most complex system considered in this work. Despite its nonlinear and chaotic behaviour, the neural network achieved high prediction accuracy when training and evaluation were performed under the same physical conditions. However, the chaos generalisation experiments demonstrated an important limitation. A network trained on one trajectory did not necessarily generalise well to trajectories generated from nearby initial conditions once the divergence between the trajectories became sufficiently large. This behaviour reflects the intrinsic sensitivity to initial conditions of chaotic systems \cite{lorenz1963deterministic,strogatz2015nonlinear} rather than an inability of the neural network to fit an individual trajectory.

Overall, the experiments suggest that the primary factor determining prediction difficulty is not whether a system is linear or nonlinear, but rather the complexity of the trajectory that must be approximated. Systems exhibiting higher oscillation frequencies, longer prediction horizons, or chaotic divergence generally require the network to approximate more rapidly varying functions \cite{spectralbias2024}, leading to increased prediction errors even when the underlying governing equations remain deterministic.

Several limitations should also be acknowledged. Only a single feedforward MLP architecture was considered, and all experiments focused on deterministic simulations without measurement noise. Future work could investigate recurrent neural networks, transformer-based models, or neural operators capable of exploiting temporal dependencies more effectively \cite{hochreiter1997long,cho2014learning}. Additionally, extending the analysis to longer prediction horizons, more strongly chaotic regimes, and physics-informed neural networks could provide further insight \cite{raissi2019physics} into the relationship between machine learning models and nonlinear dynamical systems.

In summary, this work demonstrates that relatively simple neural networks can accurately approximate a broad range of physical systems, including nonlinear and chaotic dynamics, provided that the trajectories remain sufficiently smooth and the prediction task stays within the interpolation regime. The results highlight both the potential and the limitations of purely data-driven approaches for modelling physical systems and provide a systematic comparison of how different dynamical properties influence neural network performance.
